% Бенчмарк UA-PAR. Кількості з paper/results/corpus_summary.json і маніфесту; правила запитів з paper/guidelines/queries.md.
\section{Бенчмарк UA-PAR}\label{sec:uapar}

\subsection{Джерела і корпус}
\uapar{} відтворює архів дослідника на матеріалі з відкритими ліцензіями українською і російською мовами (ліцензії --- у розділі~\ref{sec:availability}). Корпус має 7\,070 файлів у шести частинах. \emph{Нотатки}: 2\,945 статей Вікіпедії (1\,960 українських, 985 російських) з 25 тем, переважно технічних, серед них антени, радіо, LTE, програмно-визначене радіо, обробка сигналів, мікроконтролери, криптографія, машинне навчання, фізика, історія й література; кожна стаття збережена як нотатка Markdown з рядком джерела. \emph{Документи}: ще 400 статей (260 українських, 140 російських), зверстаних у 133 PDF, 134 DOCX і 133 PPTX, щоб архів мав офісні формати. \emph{Таблиці}: 150 наборів даних data.gov.ua, 100 збережено як CSV і 50 перетворено на XLSX. \emph{Код}: 1\,775 файлів з 30 репозиторіїв GitHub на українську тематику, проіндексованих як колекція коду. \emph{Зображення}: 1\,000 зображень Crossmodal-3600. \emph{Мовлення}: 800 записів FLEURS, 400 українських і 400 російських. Статті зібрано тематичним пошуком збирачем із посторінковою вибіркою і відновленням; для теми SDR не вистачило довгих статей, тож нотаток на 55 менше від запланованих 3\,000, а не доповнено слабко пов'язаними статтями.

Не кожен зібраний файл індексується: archive-index пропускає згенеровані й сторонні файли в репозиторіях, тож повний індекс EG2 містить 6\,054 документи (759 з частини коду) і 96\,908 фрагментів. Суто текстовий індекс E5 містить 4\,137 текстових документів того самого корпусу і жодних зображень чи мовлення; у розділі~\ref{sec:results} показано, що обмеження EG2 тими самими 4\,137 документами змінює його nDCG@10 на $-0{,}0015$.

\subsection{Запити}
Є 200 запитів п'яти категорій із фіксованими квотами: 40 запитів \emph{на відомий документ}, що описують конкретний документ, не цитуючи його назви; 60 \emph{тематичних} запитів, на які можуть відповідати кілька документів; 40 \emph{технічних} запитів з ідентифікаторами, номерами моделей чи стандартами в тому вигляді, як їх пишуть люди; 30 \emph{міжмовних} запитів російською до української цілі або навпаки; і 30 запитів \emph{без відповіді} --- правдоподібних питань про те, чого в корпусі немає, часто з вигаданими номерами моделей. 126 запитів українською і 74 російською. Кожен із 170 запитів з відповіддю написано для однієї цільової нотатки; вона записана, але не є єдиним релевантним документом. Валідатор перевіряє квоти й частки мов, існування кожної цілі, те, що жоден тематичний чи міжмовний запит не має з ціллю спільної послідовності понад п'ять токенів поспіль (жодного копіювання), і те, що кожен запит без відповіді не має жодного лексичного збігу в індексі. Усі 200 запитів проходять перевірку.

Запити написала GPT-6-Luna через Codex CLI за цими правилами. Оскільки автор запитів не має оцінювати власні запити, оцінки GPT ніколи не входять до міток релевантності; вони використовуються лише для порівняння оцінювання, яке не відокремлює самопреференцію від загальної поблажливості (розділ~\ref{sec:results}). Запити без відповіді за побудовою не мають лексичних збігів, тож BM25 нічого на них не повертає; порівняння систем на цих запитах --- це перевірка, чи щільна система колись утримується від відповіді, а не чесне змагання лексичного і щільного пошуку.

\subsection{Оцінювання релевантності}
Пул оцінювання — об'єднання перших десяти документів кожної з десяти систем: 5\,572 пари «запит--документ» для 200 запитів. Окремо перемішані аркуші приховують назви пошукових систем. Суддя бачить запит, назву й початок документа, разом до 1\,200 символів; для аудіозаписів додано транскрипт джерела. Шкала: 0 — не стосується запиту, 1 — частково релевантний, 2 — прямо відповідає на запит. Це оцінки видимого фрагмента, а не результат перегляду повного документа.

Початковий основний суддя — Gemini~3.8 Flash через \code{agy}, пакетами по 50 пар. Gemma~4 12B через Ollama дає основні оцінки для перевірки стійкості висновків; Gemma~4 E4B — додаткові оцінки з відкритими вагами. Запити склала GPT-6-Luna; її оцінки наведено лише для порівняння, до основних і контрольних міток релевантності вони не входять. Усі початкові аркуші збережено. Gemini та Gemma мають спільного постачальника, тому окреме оцінювання не означає незалежності навчальних даних або помилок.

Початкові Gemini та Gemma~4 12B розійшлися на 1\,091 парі. Для варіанта після перегляду окремий прохід Gemini переоцінює ці пари, бачачи обидві початкові оцінки під нейтральними позначеннями A/B зі збалансованим порядком подання. Узгоджені оцінки зберігаються. Це перегляд тією самою моделлю, а не незалежний арбітраж третього судді. Результати пошуку порівнюються за початковими Gemini, початковими Gemma~4 12B та Gemini після перегляду. Основного суддю обрано після перегляду узгодженості й діагностики цільових документів та запитів без відповіді, до поданого сукупного порівняння ефективності систем; попередньої реєстрації цього вибору не було. Завершена перевірка релевантності людиною не заявляється.

Усі 170 запитів із відповіддю мають цільові документи з нотаток; до пулу глибини 10 потрапили 134 цілі. Отже, знаменник діагностики оцінок цільових документів — 134, а не весь набір цілей. Тридцять запитів, позначених як такі, що не мають відповіді, аналізуються окремо.

\subsection{Мультимодальні підмножини}
Пошук зображень і мовлення за текстом оцінюється поза пулом бенчмарку за парами «підпис--зображення» та «транскрипт--запис» із вихідних наборів, без окремо зібраних оцінок релевантності. Для зображень запит --- перший український чи російський підпис зображення в Crossmodal-3600, а правильна відповідь --- це зображення серед усіх 1\,000. Для мовлення запит --- розшифровка речення FLEURS, а правильні відповіді --- записи цього речення; FLEURS n-сторонньо паралельний, тож російська розшифровка може шукати українські записи того самого речення і навпаки. Кожну умову оцінено трьома способами: за аудіовектором запису, за його розшифровкою, закодованою як текст, і за кращим із двох. Коли запит --- розшифровка того самого запису, текстова й комбінована оцінки є верхньою межею, а не результатом пошуку; твердження спираються на умови лише з аудіо і на міжмовні умови.
